% !TeX program = xelatex
\documentclass[11pt,a4paper]{article}

% ---------- 宏包 ----------
\usepackage[UTF8]{ctex}
\usepackage{geometry}
\geometry{margin=2.2cm}
\usepackage{amsmath,amssymb,amsthm,mathtools}
\usepackage{bm}
\usepackage{enumitem}
\usepackage{hyperref}
\hypersetup{colorlinks=true,linkcolor=blue,urlcolor=blue,citecolor=blue}

% ---------- 常用符号与宏 ----------
\newcommand{\R}{\mathbb{R}}
\newcommand{\C}{\mathbb{C}}
\newcommand{\gl}{\mathfrak{gl}}
\newcommand{\so}{\mathfrak{so}}
\newcommand{\oo}{\mathfrak{o}}
\newcommand{\Lie}{\mathrm{Lie}}
\newcommand{\X}{\mathfrak{X}} % 向量场空间
\newcommand{\dd}{\,\mathrm{d}}
\newcommand{\id}{\mathrm{id}}
\newcommand{\Te}{T_e}
\newcommand{\Ad}{\mathrm{Ad}}

% ---------- 定理环境 ----------
\theoremstyle{definition}
\newtheorem{definition}{定义}[section]
\newtheorem{example}[definition]{例}
\newtheorem{remark}[definition]{注}

\theoremstyle{plain}
\newtheorem{theorem}[definition]{定理}
\newtheorem{proposition}[definition]{命题}
\newtheorem{lemma}[definition]{引理}
\newtheorem{corollary}[definition]{推论}

% ---------- 标题 ----------
\title{向量场、李括号与李群的李代数}
\author{}
\date{}

\begin{document}
\maketitle

\tableofcontents

\section{向量场与映射的关系：$F$-related、推前（pushforward）}

\begin{definition}[ $F$-related 向量场]
设 $F:M\to N$ 为光滑映射，$X\in\X(M)$，$Y\in\X(N)$。
称 $Y$ 与 $X$ \textbf{$F$-相关}（$F$-related），若对一切 $p\in M$ 有
\[
\dd F_p(X_p)=Y_{F(p)}.
\]
\end{definition}

\begin{proposition}[微分同胚下的推前唯一存在性（命题 8.19）]
设 $M,N$ 为光滑流形（可带边界），$F:M\to N$ 为微分同胚。则对每个 $X\in\X(M)$，存在唯一的 $Y\in\X(N)$ 与 $X$ $F$-相关。
并且该向量场可显式给出为
\[
(F_*X)_q=\dd F_{F^{-1}(q)}\bigl(X_{F^{-1}(q)}\bigr),\qquad q\in N.
\]
\end{proposition}

\begin{definition}[推前（pushforward）]
在命题 8.19 的情形中，记唯一的 $F$-相关向量场为 $F_*X$，
称其为 $X$ 在 $F$ 下的\textbf{推前}（pushforward）。
\end{definition}

\begin{example}[计算推前（例 8.20）]
令 $M,N\subset \R^2$，定义微分同胚 $F:M\to N$，并给定 $X\in\X(M)$。
可用上式
\[
(F_*X)_q=\dd F_{F^{-1}(q)}(X_{F^{-1}(q)})
\]
结合 $F^{-1}$ 与雅可比矩阵直接计算 $F_*X$ 的坐标表达式。
\end{example}

\begin{corollary}[推前与对函数的作用（推论 8.21）]
设 $F:M\to N$ 为微分同胚，$X\in\X(M)$。对任意 $f\in C^\infty(N)$，
\[
\bigl((F_*X)f\bigr)\circ F = X(f\circ F).
\]
\end{corollary}

\section{子流形上的切向量场与限制}

\begin{definition}[向量场在子流形上的切向性]
设 $S\subseteq M$ 为浸入或嵌入子流形（可带边界）。向量场 $X\in\X(M)$
称为在 $p\in S$ \textbf{切于 $S$}，若 $X_p\in T_pS\subseteq T_pM$；
称 $X$ \textbf{切于 $S$}，若对一切 $p\in S$ 都切于 $S$。
\end{definition}

\begin{proposition}[切向性的函数刻画（命题 8.22）]
设 $M$ 为光滑流形，$S\subseteq M$ 为嵌入子流形（可带边界），$X\in\X(M)$。
则 $X$ 切于 $S$ 当且仅当：对每个满足 $f|_S=0$ 的 $f\in C^\infty(M)$，都有
\[
(Xf)|_S=0.
\]
\end{proposition}

\begin{proposition}[向量场限制到子流形（命题 8.23）]
设 $M$ 为光滑流形，$S\subseteq M$ 为浸入子流形（可带边界），
$\iota:S\hookrightarrow M$ 为包含映射。若 $Y\in\X(M)$ 切于 $S$，
则存在唯一的光滑向量场 $X\in\X(S)$，记为 $Y|_S$，
使得 $X$ 与 $Y$ \textbf{$\iota$-相关}（即 $\dd \iota(X_p)=Y_p$ 对 $p\in S$ 成立）。
\end{proposition}

\section{李括号：定义、例子与坐标公式}

\begin{definition}[复合算子与李括号]
设 $X,Y\in\X(M)$。对 $f\in C^\infty(M)$，定义
\[
[X,Y]f := X(Yf)-Y(Xf).
\]
由此得到一个新向量场 $[X,Y]\in\X(M)$，称为 $X,Y$ 的\textbf{李括号}（Lie bracket）。
\end{definition}

\begin{example}[复合一般不满足乘法法则（例 8.24）]
在 $\R^2$ 上令 $X=\frac{\partial}{\partial x}$，$Y=x\frac{\partial}{\partial y}$。
则算子 $f\mapsto X(Yf)$ 一般不是导子（不满足 Leibniz 法则），
但其交换子 $[X,Y]=XY-YX$ 是向量场。
\end{example}

\begin{lemma}[李括号是光滑向量场（引理 8.25）]
任意光滑流形 $M$ 上，任意 $X,Y\in\X(M)$ 的李括号 $[X,Y]$ 仍是一个光滑向量场。
\end{lemma}

\begin{proposition}[李括号的坐标表达（命题 8.26）]
设在局部坐标 $(x^i)$ 下
\[
X=X^i\frac{\partial}{\partial x^i},\qquad
Y=Y^j\frac{\partial}{\partial x^j},
\]
则
\[
[X,Y]=\Bigl(X^i\frac{\partial Y^j}{\partial x^i}-Y^i\frac{\partial X^j}{\partial x^i}\Bigr)\frac{\partial}{\partial x^j}.
\]
也可写为
\[
[X,Y]=(XY^j-YX^j)\frac{\partial}{\partial x^j}.
\]
\end{proposition}

\begin{proposition}[坐标向量场的李括号（式 8.10）]
在任意光滑坐标图中，
\[
\Bigl[\frac{\partial}{\partial x^i},\frac{\partial}{\partial x^j}\Bigr]=0\qquad(\forall\, i,j).
\]
\end{proposition}

\begin{example}[用坐标公式计算（例 8.27）]
在 $\R^3$ 上取
\[
X= x\frac{\partial}{\partial x}+\frac{\partial}{\partial y}+x(y+1)\frac{\partial}{\partial z},\qquad
Y=\frac{\partial}{\partial x}+y\frac{\partial}{\partial z},
\]
则可用命题 8.26 的坐标公式计算 $[X,Y]$。
\end{example}

\begin{proposition}[李括号的基本性质（命题 8.28）]
对任意 $X,Y,Z\in\X(M)$ 与 $a,b\in\R$：
\begin{enumerate}[label=(\alph*)]
\item \textbf{双线性：}\quad $[aX+bY,Z]=a[X,Z]+b[Y,Z]$，且 $[Z,aX+bY]=a[Z,X]+b[Z,Y]$。
\item \textbf{反对称：}\quad $[X,Y]=-[Y,X]$。
\item \textbf{Jacobi 恒等式：}\quad $[X,[Y,Z]]+[Y,[Z,X]]+[Z,[X,Y]]=0$。
\item \textbf{对函数的乘法法则：}\quad 对 $f,g\in C^\infty(M)$，
\[
[fX,gY]=fg[X,Y]+(fXg)Y-(gYf)X. \tag{8.11}
\]
\end{enumerate}
\end{proposition}

\section{自然性（naturality）与推论：李括号与映射相容}

\begin{proposition}[李括号的自然性（命题 8.30）]
设 $F:M\to N$ 为光滑映射，且 $X_1,X_2\in\X(M)$、$Y_1,Y_2\in\X(N)$ 满足
$X_i$ 与 $Y_i$ $F$-相关（$i=1,2$）。则
\[
[X_1,X_2]\ \text{与}\ [Y_1,Y_2]\ \text{$F$-相关}.
\]
\end{proposition}

\begin{corollary}[李括号的推前（推论 8.31）]
若 $F:M\to N$ 为微分同胚，$X_1,X_2\in\X(M)$，则
\[
F_*[X_1,X_2]=[F_*X_1,F_*X_2].
\]
\end{corollary}

\begin{corollary}[切向向量场的括号仍切向（推论 8.32）]
设 $S\subseteq M$ 为浸入子流形（可带边界），若 $Y_1,Y_2\in\X(M)$ 都切于 $S$，
则 $[Y_1,Y_2]$ 也切于 $S$。
\end{corollary}

\section{李群与左不变向量场：$\Lie(G)$ 与李代数}

\begin{definition}[左平移与左不变向量场]
设 $G$ 为李群。对 $g\in G$，左平移 $L_g:G\to G$ 定义为 $L_g(h)=gh$。
向量场 $X\in\X(G)$ 称为\textbf{左不变}，若对每个 $g\in G$，
$X$ 与自身在 $L_g$ 下相关，即
\[
\dd(L_g)_{g'}(X_{g'})=X_{gg'}\qquad(\forall\, g,g'\in G).
\]
等价地（因 $L_g$ 为微分同胚）：
\[
(L_g)_*X=X\qquad(\forall\, g\in G). \tag{8.12}
\]
\end{definition}

\begin{proposition}[左不变向量场对李括号封闭（命题 8.33）]
设 $G$ 为李群，$X,Y$ 为 $G$ 上的光滑左不变向量场，则 $[X,Y]$ 仍是左不变向量场。
\end{proposition}

\begin{definition}[李群的李代数 $\Lie(G)$]
记 $\Lie(G)$ 为李群 $G$ 上\textbf{所有光滑左不变向量场}的集合。
在通常的向量场加法与数乘下，$\Lie(G)$ 是向量空间；
在李括号 $[\cdot,\cdot]$ 下，$\Lie(G)$ 构成一个李代数。
\end{definition}

\begin{example}[李代数的典型来源（例 8.36）]
\begin{enumerate}[label=(\alph*)]
\item $\X(M)$ 在李括号下构成李代数。
\item 若 $G$ 为李群，则所有左不变向量场组成 $\Lie(G)$，它是 $\X(G)$ 的李子代数。
\item $M(n,\R)$ 在交换子括号 $[A,B]=AB-BA$ 下成为李代数，记为 $\gl(n,\R)$。
\item 类似地，$M(n,\C)$ 在交换子括号下为实李代数，记为 $\gl(n,\C)$（实维 $2n^2$）。
\item 若 $V$ 为向量空间，则 $\gl(V)=\mathrm{End}(V)$ 在 $[A,B]=A\circ B-B\circ A$ 下为李代数。
\item 若任意向量空间上定义括号恒为零，则得到\textbf{阿贝尔}李代数。
\end{enumerate}
\end{example}

\begin{theorem}[取值映射同构（定理 8.37）]
设 $G$ 为李群。取值映射
\[
\varepsilon:\Lie(G)\to T_eG,\qquad \varepsilon(X)=X_e,
\]
是向量空间同构。因此 $\Lie(G)$ 有限维，且
\[
\dim \Lie(G)=\dim G.
\]
\end{theorem}

\begin{definition}[由 $v\in T_eG$ 生成的左不变向量场（式 8.13）]
对任意 $v\in T_eG$，定义（粗略）向量场 $v^L$：
\[
(v^L)_g=\dd(L_g)_e(v),\qquad g\in G. \tag{8.13}
\]
该向量场实际上是光滑左不变向量场，并满足 $(v^L)_e=v$。
\end{definition}

\begin{corollary}[左不变粗略向量场必光滑（推论 8.38）]
李群 $G$ 上每个左不变的粗略向量场都是光滑的。
\end{corollary}

\begin{corollary}[李群可平行化（推论 8.39）]
每个李群都存在一组左不变的光滑\textbf{全局标架}（global frame），因此每个李群都是可平行化的（parallelizable）。
\end{corollary}

\begin{example}[熟悉李群的李代数（例 8.40）]
\begin{enumerate}[label=(\alph*)]
\item $\R^n$（加法群）：左不变向量场恰为常系数向量场，故 $\Lie(\R^n)\cong \R^n$（阿贝尔）。
\item 圆周群 $S^1$：$\Lie(S^1)\cong \R$（阿贝尔）。
\item $n$-维环面 $T^n=(S^1)^n$：$\Lie(T^n)\cong \R^n$（阿贝尔）。
\item 一般线性群 $GL(n,\R)$：其李代数与 $\gl(n,\R)$ 自然同构（见后述命题 8.41）。
\end{enumerate}
\end{example}

\section{$GL(n,\R)$ 与矩阵李代数：$\Lie(GL(n,\R))\cong\gl(n,\R)$}

\begin{proposition}[一般线性群的李代数（命题 8.41）]
自然映射的复合
\[
\Lie\bigl(GL(n,\R)\bigr)\ \longrightarrow\ T_{I_n}GL(n,\R)\ \longrightarrow\ \gl(n,\R) \tag{8.14}
\]
给出 $\Lie(GL(n,\R))$ 与矩阵李代数 $\gl(n,\R)$ 之间的\textbf{李代数同构}。
\end{proposition}

\begin{corollary}[抽象向量空间版本（推论 8.42）]
设 $V$ 为有限维实向量空间，则自然同构的复合
\[
\Lie\bigl(GL(V)\bigr)\ \longrightarrow\ T_{\id}GL(V)\ \longrightarrow\ \gl(V) \tag{8.16}
\]
给出 $\Lie(GL(V))$ 与 $\gl(V)$ 的李代数同构。
\end{corollary}

\section{李群同态诱导的李代数同态}

\begin{theorem}[诱导李代数同态（定理 8.44）]
设 $G,H$ 为李群，$\mathfrak g=\Lie(G)$，$\mathfrak h=\Lie(H)$。
若 $F:G\to H$ 为李群同态，则对每个 $X\in\mathfrak g$，存在唯一 $Y\in\mathfrak h$
与 $X$ $F$-相关。记该 $Y$ 为 $F_*X$，则由此定义的映射
\[
F_*:\mathfrak g\to\mathfrak h
\]
是一个\textbf{李代数同态}。
\end{theorem}

\begin{proposition}[诱导同态的性质（命题 8.45）]
\begin{enumerate}[label=(\alph*)]
\item 恒等同态诱导恒等：$(\id_G)_*=\id_{\Lie(G)}$。
\item 复合相容：若 $F_1:G\to H$、$F_2:H\to K$ 为李群同态，则
\[
(F_2\circ F_1)_*=(F_2)_*\circ(F_1)_*:\Lie(G)\to\Lie(K).
\]
\item 同构推出李代数同构：若 $F:G\to H$ 为李群同构，则 $F_*:\Lie(G)\to\Lie(H)$ 为李代数同构。
\end{enumerate}
\end{proposition}

\section{李子群的李代数}

\begin{theorem}[李子群的李代数（定理 8.46）]
设 $H\subseteq G$ 为李子群，$\iota:H\hookrightarrow G$ 为包含映射。
则存在李子代数 $\mathfrak h\subseteq \Lie(G)$ 与 $\Lie(H)$ 典范同构，且
\[
\mathfrak h=\iota_*\bigl(\Lie(H)\bigr)
=\{\,X\in \Lie(G): X_e\in T_eH\,\}. \tag{8.17}
\]
\end{theorem}

\begin{example}[正交群的李代数（例 8.47）]
$O(n)$ 是 $GL(n,\R)$ 的李子群，其在单位元处的切空间满足
\[
T_{I_n}O(n)=\{\,B\in\gl(n,\R):B^T+B=0\,\},
\]
即所有 $n\times n$ 反对称矩阵。该子空间记为 $\oo(n)$，
并与 $\Lie(O(n))$ 典范同构，且 $\oo(n)$ 是 $\gl(n,\R)$ 的李子代数。
\end{example}

\section{$GL(n,\C)$ 的李代数}

\begin{proposition}[$GL(n,\C)$ 的李代数（命题 8.48）]
由自然同构
\[
\Lie(GL(n,\C))\xrightarrow{\ \varepsilon\ }T_{I_n}GL(n,\C)\xrightarrow{\ \varphi\ }\gl(n,\C) \tag{8.18}
\]
的复合得到 $\Lie(GL(n,\C))$ 与矩阵李代数 $\gl(n,\C)$ 之间的\textbf{李代数同构}
（将其视为\textbf{实}李代数）。
\end{proposition}

\section{李代数表示、Ado 定理与矩阵化推论}

\begin{definition}[李代数表示与忠实性]
设 $\mathfrak g$ 为有限维李代数。称李代数同态
\[
\varphi:\mathfrak g\to\gl(V)
\]
为 $\mathfrak g$ 的（有限维）\textbf{表示}。若 $\varphi$ 为单射，则称该表示\textbf{忠实}（faithful）。
\end{definition}

\begin{theorem}[Ado 定理（定理 8.49）]
每个有限维实李代数都存在一个忠实的有限维表示。
\end{theorem}

\begin{corollary}[有限维李代数可嵌入矩阵李代数（推论 8.50）]
每个有限维实李代数都同构于某个矩阵李代数 $\gl(n,\R)$（带交换子括号）中的一个李子代数。
\end{corollary}

\end{document}
